\chapter{CONCLUSION AND FUTURE SCOPE}

\section{Summary of the Project}

This report presented \textbf{CodeAudit AI}, a novel multi-agent Retrieval-Augmented Generation (RAG) framework designed to authenticate technical skill claims on candidate resumes against publicly available source code repositories on GitHub \cite{ref20, ref38}. Traditional resume screening platforms and Applicant Tracking Systems (ATS) rely almost exclusively on keyword matching and candidate self-reporting, creating severe vulnerabilities to skill inflation, exaggerated experience, and AI-generated resume stuffing \cite{ref1, ref9, ref17}. CodeAudit AI bridges this critical verification gap by anchoring technical skill evaluation in concrete, functional code artifacts \cite{ref12, ref7, ref35}.

The system architecture decomposes the verification workflow into four specialized autonomous agents operating in sequence:
\begin{enumerate}[leftmargin=*, label=\arabic*.]
    \item \textbf{Claim Extraction Agent:} Parses uploaded candidate PDF resumes using PyMuPDF (fitz) and applies regular expression scanning with word-boundary anchors to extract top-5 technical skill claims \cite{ref40}.
    \item \textbf{Code Ingestion Agent:} Asynchronously queries the GitHub REST API v3 using \texttt{aiohttp} to retrieve candidate public repositories, applying file extension whitelists and a 100KB per-file cap to filter binary blobs \cite{ref12, ref35}.
    \item \textbf{RAG Indexing Agent:} Segments ingested source files into character-level chunks (800 chars, 100 overlap), embeds chunks into 384-dimensional dense vectors using \texttt{all-MiniLM-L6-v2}, and stores embeddings in ChromaDB with mandatory candidate-level metadata filtering \cite{ref31, ref33, ref41}.
    \item \textbf{LLM Judge Agent:} Orchestrated via LangChain, retrieves top-5 semantically relevant code chunks from ChromaDB and applies Chain-of-Thought (CoT) reasoning to emit deterministic binary verdicts (\texttt{Verified} or \texttt{Unsubstantiated}) with explicit code citations \cite{ref10, ref27, ref25, ref42}.
\end{enumerate}

\section{Key Contributions}

The primary technical and practical contributions of this project are summarized as follows:

\begin{enumerate}[leftmargin=*, label=\arabic*.]
    \item \textbf{Novel Multi-Agent Verification Architecture:} Formulated a 4-phase autonomous agent pipeline for automated codebase claim auditing, isolating responsibilities across extraction, ingestion, indexing, and judgment \cite{ref37, ref22, ref43}.
    \item \textbf{Empirical Evidence Grounding:} Integrated RAG-based vector search over real-world candidate source code, eliminating LLM hallucination in technical evaluation by requiring explicit code syntax citations \cite{ref30, ref31, ref32}.
    \item \textbf{Strict Multi-Tenant Isolation:} Designed candidate-level metadata filtering in ChromaDB, guaranteeing complete data segregation in multi-candidate screening environments \cite{ref33, ref34}.
    \item \textbf{Rigorous Benchmark Dataset:} Constructed a 30-case benchmark across 5 technical domains, incorporating 20 True Positives, 10 True Negatives, and 2 adversarial edge cases to evaluate system robustness \cite{ref20, ref8}.
    \item \textbf{Zero False Positive Rate:} Demonstrated that strict indicator-based prompt tuning (System B) achieves 100.0\% Precision and 100.0\% Specificity, ensuring that no unverified claim is validated \cite{ref36}.
\end{enumerate}

\section{Key Findings and Empirical Results}

Empirical evaluation on the 30-case benchmark dataset yielded compelling results across system variants:
\begin{itemize}[leftmargin=*]
    \item \textbf{Baseline (No RAG):} Achieved 60.0\% accuracy, 70.0\% precision, and 40.0\% specificity, demonstrating the severe limitations of naive keyword matching.
    \item \textbf{System A (RAG + Moderate):} Reached 93.3\% accuracy, 90.9\% precision, 100.0\% recall, and 80.0\% specificity, showing significant improvement but occasionally validating superficial mentions.
    \item \textbf{System B (RAG + Strict):} Achieved 93.3\% accuracy, 100.0\% Precision, 90.0\% Recall, a 94.7\% F1-score, and 100.0\% Specificity, completely eliminating false positives.
\end{itemize}

Category-wise evaluation confirmed 100.0\% accuracy and F1-score across Databases (DATA), Software Engineering (SWE), and DevOps/Tools (TOOL) categories, with minor dips in Machine Learning (88.9\% accuracy) and Web Development (83.3\% accuracy) due to abstract library wrappers (TC-28) and configuration-only framework declarations (TC-29). End-to-end processing latency averaged 15 to 45 seconds per candidate portfolio, confirming real-world deployment viability \cite{ref41}.

\section{Limitations of the Current System}

Despite its strong performance, CodeAudit AI exhibits specific technical limitations:
\begin{enumerate}[leftmargin=*, label=\arabic*.]
    \item \textbf{Public Repository Scope:} Current ingestion is limited to public GitHub repositories, excluding private enterprise codebases and alternate platforms (GitLab, Bitbucket).
    \item \textbf{File Size and Type Caps:} Source code ingestion is capped at 100KB per file, excluding compiled binaries, heavy data notebooks, and non-whitelisted extensions.
    \item \textbf{Wrapper and Abstraction Blindspots:} Indirect library imports wrapped inside custom modules can trigger conservative false negatives under strict prompt rules.
    \item \textbf{API Dependency and Latency:} Processing speed depends on external GitHub API rate limits and network bandwidth during asynchronous ingestion \cite{ref12}.
\end{enumerate}

\section{Future Research Scope}

Future extensions of CodeAudit AI include the following research directions:
\begin{enumerate}[leftmargin=*, label=\arabic*.]
    \item \textbf{Multi-Platform \& Private Repo Integration:} Extending ingestion to GitLab, Bitbucket, and OAuth-authenticated private enterprise repositories \cite{ref12, ref15}.
    \item \textbf{Advanced AST \& Data-Flow Parsing:} Incorporating tree-sitter static analysis to parse Abstract Syntax Trees (ASTs) alongside vector search \cite{ref67, ref8}.
    \item \textbf{Fine-Tuned Local Code LLMs:} Deploying fine-tuned Code Llama or StarCoder instances locally via Ollama to eliminate external API costs \cite{ref70, ref17}.
    \item \textbf{Integration with Applicant Tracking Systems (ATS):} Developing REST API webhooks for seamless integration into enterprise ATS platforms (Greenhouse, Lever) \cite{ref1, ref4}.
\end{enumerate}

\section{Social and Ethical Implications}

CodeAudit AI directly addresses key ethical mandates in algorithmic recruitment \cite{ref4, ref14, ref22, ref65, ref66}. By substituting subjective recruiter bias and keyword-stuffing exploits with transparent, evidence-grounded code citations, the system promotes fair, meritocratic technical evaluation \cite{ref22, ref36}. Strict metadata isolation ensures compliance with international data privacy regulations (EU AI Act, UNESCO Recommendations), establishing an auditable foundation for future AI hiring technology \cite{ref65, ref66}.



\section{Broader Impact and Industry Deployment Roadmap}

The successful implementation and empirical validation of CodeAudit AI establish a transformative blueprint for modern technical recruitment \cite{ref1, ref20, ref23}. Beyond initial resume screening, the multi-agent RAG paradigm introduced in this research has broader implications for automated credential verification, technical talent analytics, and enterprise talent acquisition workflows \cite{ref4, ref7, ref22}.

\subsection{Enterprise System Integration}

For enterprise organizations processing thousands of technical applications monthly, CodeAudit AI can be integrated directly into existing Human Resource Information Systems (HRIS) and Applicant Tracking Systems (ATS) through asynchronous REST API webhooks \cite{ref1, ref12}. When a candidate submits an application via platforms like Greenhouse, Lever, or Workday, an automated webhook trigger initiates the 4-phase audit pipeline in the background. Audit results—comprising verified skill badges, ChromaDB distance metrics, and source code citations—are rendered directly within the recruiter's candidate dashboard prior to scheduling technical interviews \cite{ref41, ref44}.

\subsection{Continuous Talent Telemetry}

Future enterprise deployments can extend CodeAudit AI to provide continuous talent telemetry across candidate career trajectories \cite{ref15, ref24}. By periodically indexing updated public GitHub repositories, technical recruiters can track candidate skill acquisition, framework adoption, and open-source contribution patterns over time. This continuous evaluation model replaces static point-in-time resume screening with dynamic, empirical talent analytics, ensuring that organizational hiring pipelines remain aligned with rapidly evolving software engineering ecosystems \cite{ref6, ref25, ref36}.

\section{Concluding Remarks}

CodeAudit AI demonstrates that combining multi-agent LLM orchestration with Retrieval-Augmented Generation provides a highly effective solution to the long-standing problem of technical resume fraud. By grounding skill evaluation in verifiable source code evidence, the system achieves perfect precision and specificity, delivering a reliable, fair, and scalable tool for modern technical hiring.